% =============================================================================
% Astro-Amriit — Illustrated Guide (v0.9)
% Build: node media/pdf/shots.mjs  (screenshots + figure files in work/)
%        cd media/pdf && latexmk guide.tex   (PDF in out/)
% Every screenshot comes from the shipped app (dist/Astro-Amriit v0.9.html)
% with the public example chart used throughout the user guide.
% =============================================================================
\documentclass[11pt, a4paper]{article}

% --- Geometry, fonts, microtypography ----------------------------------------
\usepackage[margin=2.2cm, headheight=15pt]{geometry}
\usepackage[T1]{fontenc}
\usepackage{lmodern}
\usepackage[utf8]{inputenc}
% arc minutes and seconds, as the app writes them (14°53′20″)
\DeclareUnicodeCharacter{2032}{\ensuremath{'}}
\DeclareUnicodeCharacter{2033}{\ensuremath{''}}
\usepackage{microtype}
\usepackage{parskip}
\usepackage{textcomp}
\usepackage{amsmath, amssymb}

% --- Tables, figures, lists, boxes --------------------------------------------
\usepackage{booktabs}
\usepackage{tabularx}
\usepackage{array}
\usepackage{longtable}
\usepackage{graphicx}
\graphicspath{{work/}}
\usepackage{caption}
\usepackage[section]{placeins}   % figures stay inside their own section
\usepackage{enumitem}
\usepackage{listings}
\usepackage{mdframed}
\usepackage{xcolor}
\usepackage{tikz}
\usetikzlibrary{positioning}
\usepackage{titlesec}
\usepackage{fancyhdr}
\usepackage[htt]{hyphenat}
\usepackage{etoolbox}

% --- House colours + MIDSEA accent (headings, boxes, tags) --------------------
\definecolor{headerblue}{RGB}{0,70,127}
\definecolor{noteblue}{RGB}{220,235,245}
\definecolor{warnyellow}{RGB}{255,248,220}
\definecolor{codegrey}{RGB}{245,245,245}
% =============================================================================
% MIDSEA colourful accent preamble fragment
% -----------------------------------------------------------------------------
% The colourful report style used for explainers (MIDSEA2026 reports).
% \input this AFTER \usepackage{xcolor} and the base house colours, and BEFORE
% \usepackage{hyperref}. Requires (loaded in the main preamble): xcolor,
% mdframed, titlesec, fancyhdr, listings. Safe to \input over the plain house
% style: it only ADDS colours/boxes/tags and recolours headings + warnbox.
% =============================================================================

% --- Base house colours (redefined harmlessly if already present) ------------
\definecolor{headerblue}{RGB}{0,70,127}
\definecolor{noteblue}{RGB}{220,235,245}
\definecolor{warnyellow}{RGB}{255,248,220}
\definecolor{codegrey}{RGB}{245,245,245}

% --- MIDSEA accent palette ---------------------------------------------------
\definecolor{midseapurple}{RGB}{75,0,130}
\definecolor{midseamagenta}{RGB}{199,21,133}
\definecolor{midseaorange}{RGB}{230,126,0}
\definecolor{midseared}{RGB}{200,30,60}
\definecolor{midseagreen}{RGB}{0,140,90}
\definecolor{keygreen}{RGB}{223,242,233}   % key-idea box background
\definecolor{qpink}{RGB}{250,228,242}      % question box background
\definecolor{derivlav}{RGB}{238,232,248}   % derivation box background
\definecolor{netpurple}{HTML}{6E07C5}      % "in plain words" tag

% --- Colourful section headings (purple text, magenta rule) ------------------
\titleformat{\section}{\large\bfseries\color{midseapurple}}{\thesection}{0.6em}{}[{\color{midseamagenta}\titlerule[1.2pt]}]
\titleformat{\subsection}{\normalsize\bfseries\color{headerblue}}{\thesubsection}{0.5em}{}
\titleformat{\subsubsection}{\normalsize\itshape\color{midseaorange!85!black}}{\thesubsubsection}{0.5em}{}

% --- Coloured running heads + footrule ---------------------------------------
% (set \rhead/\lhead text yourself, keeping the \color{midseapurple} prefix)
\renewcommand{\headrulewidth}{0.4pt}
\renewcommand{\footrulewidth}{0.2pt}

% --- Callout boxes -----------------------------------------------------------
\newmdenv[backgroundcolor=noteblue, linecolor=headerblue, linewidth=1pt, roundcorner=4pt,
  innerleftmargin=10pt, innerrightmargin=10pt, innertopmargin=8pt, innerbottommargin=8pt]{notebox}
\newmdenv[backgroundcolor=warnyellow, linecolor=midseaorange, linewidth=1pt, roundcorner=4pt,
  innerleftmargin=10pt, innerrightmargin=10pt, innertopmargin=8pt, innerbottommargin=8pt]{warnbox}
\newmdenv[backgroundcolor=keygreen, linecolor=midseagreen, linewidth=1pt, roundcorner=4pt,
  innerleftmargin=10pt, innerrightmargin=10pt, innertopmargin=8pt, innerbottommargin=8pt]{keybox}
\newmdenv[backgroundcolor=qpink, linecolor=midseamagenta, linewidth=1pt, roundcorner=4pt,
  innerleftmargin=10pt, innerrightmargin=10pt, innertopmargin=8pt, innerbottommargin=8pt]{qbox}
\newmdenv[backgroundcolor=derivlav, linecolor=midseapurple, linewidth=1pt, roundcorner=4pt,
  innerleftmargin=10pt, innerrightmargin=10pt, innertopmargin=8pt, innerbottommargin=8pt]{derivbox}
\newmdenv[backgroundcolor=codegrey, linecolor=gray!40, linewidth=0.5pt, roundcorner=2pt,
  innerleftmargin=10pt, innerrightmargin=10pt, innertopmargin=6pt, innerbottommargin=6pt]{codebox}

% --- Inline coloured tags (open a box with one of these) ---------------------
\newcommand{\kidea}{{\color{midseagreen}\textbf{\textsf{Key idea.}}}\ }
\newcommand{\kpit}{{\color{midseaorange!85!black}\textbf{\textsf{Pitfall.}}}\ }
\newcommand{\kq}{{\color{midseamagenta}\textbf{\textsf{Question.}}}\ }
\newcommand{\kcode}{{\color{headerblue}\textbf{\textsf{In the code.}}}\ }
\newcommand{\ksimple}{{\color{netpurple}\textbf{\textsf{In plain words.}}}\ }

% --- Colourful code keywords -------------------------------------------------
\lstset{keywordstyle=\color{midseapurple}\bfseries, stringstyle=\color{midseagreen}}

% Usage:
%   \begin{keybox}\kidea The sentence to memorise.\end{keybox}
%   \begin{warnbox}\kpit The easy mistake / honest caveat.\end{warnbox}
%   \begin{qbox}\kq \emph{``Anticipated question,''} answered here.\end{qbox}
%   \begin{derivbox} A short derivation or why-it-holds argument.\end{derivbox}

% the app's own colours (cover, callouts)
\definecolor{plate}{HTML}{141B3D}
\definecolor{platedeep}{HTML}{0E1330}
\definecolor{plategold}{HTML}{C6A15B}
\definecolor{platedisplay}{HTML}{ECD396}
\definecolor{platemuted}{HTML}{A6AECF}

% --- Callouts on screenshots ----------------------------------------------------
\tikzset{
  callout/.style={draw=midseamagenta, line width=1.1pt, rounded corners=2pt},
  badge/.style={circle, fill=midseamagenta, text=white, font=\sffamily\bfseries\scriptsize,
                inner sep=1.2pt, minimum size=11pt, anchor=center},
}
\newcommand{\badgeinline}[1]{\tikz[baseline=(b.base)]\node[badge](b){#1};}
\newlist{calloutlist}{description}{1}
\setlist[calloutlist]{font=\normalfont, leftmargin=1.9em, labelwidth=1.4em, labelsep=0.5em,
  itemsep=1pt, topsep=4pt, before=\small\raggedright}
\newcommand{\shot}[1]{\input{work/fig/#1}}

% --- Headers, captions, layout safety ----------------------------------------
\pagestyle{fancy}
\fancyhf{}
\rhead{\small\color{midseapurple}Astro-Amriit — Illustrated Guide}
\lhead{\small\color{midseapurple}Karar Ashram Astrology Software · v0.9}
\rfoot{\small Page \thepage}
\lfoot{\small Example chart: Amitabh Bachchan, 11 Oct 1942, 16:00, Allahabad}
\captionsetup{font=small, labelfont=bf, skip=6pt}
\emergencystretch=3em
\setlength{\tabcolsep}{5pt}
\renewcommand{\arraystretch}{1.15}

% --- Hyperref last ------------------------------------------------------------
\usepackage{hyperref}
\hypersetup{colorlinks=true, linkcolor=midseapurple, urlcolor=headerblue, citecolor=midseagreen,
  pdftitle={Astro-Amriit — Illustrated Guide (v0.9)}, pdfauthor={Karar Ashram},
  pdfsubject={User guide for Astro-Amriit, the Karar Ashram astrology software}}

\newcommand{\app}{Astro-Amriit}
\newcommand{\PL}{Parashara's Light}
\newcommand{\ui}[1]{\textbf{\textsf{#1}}}   % on-screen names: buttons, tabs, labels
\newcommand{\ex}[1]{\texttt{#1}}             % something you type

\begin{document}

% =============================================================================
% Cover
% =============================================================================
\begin{titlepage}
\begin{tikzpicture}[remember picture, overlay]
  \fill[platedeep] (current page.south west) rectangle (current page.north east);
  \fill[plate] ([yshift=-3cm]current page.north west) rectangle ([yshift=7cm]current page.south east);
  % the East Indian chart plate
  \begin{scope}[shift={([yshift=-1.2cm]current page.center)}, scale=1.05, plategold, line width=1.2pt]
    \draw (-3.6,-3.6) rectangle (3.6,3.6);
    \draw (-1.2,-3.6) -- (-1.2,3.6) (1.2,-3.6) -- (1.2,3.6) (-3.6,-1.2) -- (3.6,-1.2) (-3.6,1.2) -- (3.6,1.2);
    \draw (-3.6,3.6) -- (-1.2,1.2) (3.6,3.6) -- (1.2,1.2) (-3.6,-3.6) -- (-1.2,-1.2) (3.6,-3.6) -- (1.2,-1.2);
  \end{scope}
  \node[platedisplay, font=\fontsize{46}{50}\selectfont] at ([yshift=-5.3cm]current page.north) {\app};
  \node[platemuted, font=\large\sffamily] at ([yshift=-6.6cm]current page.north) {K A R A R \ \ A S H R A M \ \ A S T R O L O G Y \ \ S O F T W A R E};
  \node[white, font=\Large] at ([yshift=-1.2cm]current page.center) {Illustrated Guide};
  \node[platemuted, font=\normalsize] at ([yshift=-1.95cm]current page.center) {every page, every panel, with the example chart};
  \node[platedisplay, font=\large] at ([yshift=4.6cm]current page.south) {Version 0.9};
  \node[platemuted, font=\small, align=center] at ([yshift=3.3cm]current page.south)
     {One HTML file \textperiodcentered{} works offline \textperiodcentered{} calculates as \PL{} 9\\[2pt]
      Screenshots taken from the shipped app, \today};
\end{tikzpicture}
\end{titlepage}

\tableofcontents
\newpage

% =============================================================================
\section{About this guide}
\label{sec:about}
% =============================================================================
\app{} is the Karar Ashram astrology software. It calculates Vedic charts the way \PL{} 9 does and draws
them in the East Indian style the family knows from Kala Chakra. This guide shows every page of the
app as it looks on screen, with numbered callouts, and explains what each part means and how to use it.

Every picture uses one public example chart: \textbf{Amitabh Bachchan, 11 October 1942, 16:00,
Allahabad (Prayagraj)}. Cast the same chart and you can compare your screen with the pictures.

\begin{keybox}\kidea \app{} is a calculating instrument. It shows the chart, the strengths, the dashas,
the divisional charts and the research tabs' rules with their sources. It does not write predictions:
the judgement is yours.\end{keybox}

\subsection{Conventions used on every page}
\begin{itemize}[itemsep=2pt]
  \item Planets are written with two letters: \ex{As} (ascendant), \ex{Su}, \ex{Mo}, \ex{Ma}, \ex{Me},
        \ex{Ju}, \ex{Ve}, \ex{Sa}, \ex{Ra}, \ex{Ke}. Each has its own colour, but the letters are always
        shown, so colour is never needed to read a page.
  \item An \ex{R} right after a planet means retrograde; a \ex{c} means combust.
  \item Signs are written with two letters: \ex{Ar Ta Ge Cn Le Vi Li Sc Sg Cp Aq Pi}.
  \item Longitudes are degrees:minutes:seconds within the sign, for example \ex{Aq 01:56:52}.
  \item Hold the mouse over a short heading (for example \ui{SB\%}) to see its full meaning.
\end{itemize}

In this guide, names you see on screen are printed \ui{like this}, and things you type
\ex{like this}. The numbered badges in each picture match the numbered list under it.

% =============================================================================
\section{Opening the app}
\label{sec:open}
% =============================================================================
The whole app is \textbf{one file} ending in \ex{.html}, for example \ex{Astro-Amriit v0.9.html}.
Copy it to your computer and double-click it: it opens in your web browser (Chrome, Edge, Firefox or
Safari). Wait for the green dot and the word \ui{ready} at the top right; this takes about a second.

It needs \textbf{no internet}. Nothing is sent anywhere and nothing is downloaded: the Swiss Ephemeris
(the planetary tables) is inside the file. When the app opens it casts a chart for \emph{this moment} at
the last place you used, so you never see an empty page.

For computers that should have the app in the Start menu or the Applications folder there is a Windows
installer and a Mac disk image with the same app. In the desktop app, \ui{Save chart} and
\ui{Save for Parashara's Light} put the file straight into your Downloads folder.

\begin{notebox}The small line at the very bottom of every page (for example \ex{Astro-Amriit v0.9.0 ·
single build …}) is the version. Quote it if you report a problem.\end{notebox}

% =============================================================================
\section{Casting a chart}
\label{sec:cast}
% =============================================================================
The chart form is at the top of the app. \autoref{fig:form} shows it filled in with the example chart.

\shot{form}

\subsection{The fields}
\begin{tabularx}{\linewidth}{>{\bfseries}l X}
\toprule
On screen & What to enter \\
\midrule
Name & Who or what the chart is for. Optional; it appears in the centre of the Rasi chart and in file names. \\
Date & The date of birth or of the question, years 1800 to 2399. \\
Time & Three boxes: hours, minutes, seconds, on the 24-hour clock (4 pm is \ex{16}). Typing two digits moves to the next box, so \ex{160037} gives 16:00:37. You may also paste \ex{16:00:37} into the hours box. \\
Place & Start typing; choose from the list to fill latitude, longitude and zone. \\
Latitude, Longitude & Decimal degrees (\ex{25.45}) or degrees, minutes, seconds (\ex{21 42 12 N}). South and west are negative. \\
Zone & Hours east of UTC for the clock time you typed; India is \ex{5.5}. \\
Daylight saving & Tick it if the clock was one hour ahead (summer time, or war time). \\
Calculation & The profile: \PL{} 9 (default), \PL{} 7, or Solar Fire / Kala Chakra. \\
\bottomrule
\end{tabularx}

\begin{warnbox}\kpit\textbf{Always read the date in words.} The date box shows the date in your
computer's order, often month/day/year (\ex{10/11/1942}). The words underneath (\ex{11 Oct 1942})
cannot be misread. Day and month swaps are the commonest mistake in hand-written birth data.\end{warnbox}

\subsection{India's war time}
From 1 September 1942 to 14 October 1945 India's clocks ran one hour ahead. When you choose an Indian
place for such a date, the app sets the zone to \ex{+6:30} by itself and shows the note
\emph{India war time was in force on this date: zone set to +6:30.} (visible in \autoref{fig:form}).
Zone \ex{6.5} without daylight saving, and zone \ex{5.5} with daylight saving ticked (how \PL{} stores it),
give the same moment. If you type latitude, longitude and zone by hand, set the zone yourself.

\subsection{The buttons}
\begin{tabularx}{\linewidth}{>{\bfseries}l X}
\toprule
Button & What it does \\
\midrule
Cast chart & Calculates the chart; the form folds away. \\
Now & Casts a chart for this very moment at the place in the form (horary, prashna). \\
Save this place & Adds the place with its coordinates and zone to your own list (up to 100). \\
Open chart file… & Opens a \PL{} \ex{.xml} chart file or an \app{} chart file and casts it. \\
Save chart & Saves an \app{} chart file (\ex{name.astro-amriit.json}). \\
Save for Parashara's Light & Saves a \PL{} chart file (\ex{name.xml}) that \PL{} can open. \\
Edit chart & After casting, opens the form again with the chart's data. \\
\bottomrule
\end{tabularx}

If a field is wrong, a red message appears under that field and nothing you typed is lost. A \PL{} file
can hold further charts for the same person (events such as ``Today''); after opening it, press
\ui{Edit chart} and they are listed under \ui{This file also holds:}.

\subsection{The chart line}
After casting, the line under the chart's name repeats the date, time, place, zone (for example
\ex{UTC+6:30}), latitude and longitude (\ex{25N27 81E51}) and the profile. Below it is the panchanga.

\shot{chartline}

% =============================================================================
\section{Page 1}
\label{sec:page1}
% =============================================================================
Page 1 follows the birth chart page of \PL{} 9, so every panel is where a \PL{} user expects it.

\shot{p1rasi}

The Navamsha (D9) beside it is drawn the same way: each box shows the sign, the house from the navamsha
ascendant, and the planets. The twelve Sarvashtakavarga numbers on the edge always add up to 337; the full
tables behind both edge numbers are on the Strengths page (\autoref{sec:strengths}).

\shot{p1planets}
\shot{p1strength}

\begin{notebox}Blank cells are blank in \PL{} too: Rahu and Ketu have no SB\%, SB\# or K7, and Ketu has no
AV.\end{notebox}

\shot{p1vim}

For the full dasha tree use the Dashas page (\autoref{sec:dashas}). Dates are in the chart's own time zone.

\shot{p1bars}
\shot{p1how}

% =============================================================================
\section{Page 1 (Experimental)}
\label{sec:exp}
% =============================================================================
Five sub-tabs gather \PL{}'s worksheets: \ui{Highlights}, \ui{Lordships \& yogas}, \ui{Aspects},
\ui{Strength worksheet} and \ui{Nandi Nadi (research)}. Each rule is written on the page where it is
used; where a \PL{} 7.0.1 printout was available, the page says it matches.

\shot{exstates}

The table spells out each code: Exalted, Moolatrikona (Rasi only), Own sign, Debilitated, Vargottama,
Dig bala, Retrograde, Combust, and ``In the 3rd / 6th / 8th / 12th'' of that chart (the 3rd in amber, as a
mild dusthana). House numbers in red on the chart mark those houses.

\shot{exconj}
\shot{exacross}
\shot{exlords}
\shot{exyogas}

\begin{notebox}Some yogas depend on the native's sex: set \ui{Native} to evaluate them. Angaraka and
Guru-Chandala are not in \PL{}'s database and are shown separately as textbook yogas. \PL{} tests yogas on
the Rasi only; choosing a divisional chart applies the same conditions to it (an extension).\end{notebox}

\shot{exmatrix}
\shot{exhouses}
\shot{exws14}

Dominance and Func.\ match a \PL{} 7.0.1 printout for all nine planets; Ishta/Kashta matches for the Sun,
Moon and Saturn and is within 1--2 virupas for the others. \PL{} 7.0.1 always leaves War blank; W* and L*
mark a textbook planetary war.

\shot{exnadi}

Nandi Nadi follows R.\,G. Rao's \emph{Bhrigu Nandi Nadi}: the Lagna is not used; each planet stands for
people and matters (Jupiter the native, Saturn work, Venus the wife, and so on). The tab also shows a card
per karaka and Jupiter's and Saturn's transits over the natal planets. It is research, not checked against
reference charts.

% =============================================================================
\section{Strengths}
\label{sec:strengths}
% =============================================================================
This page shows how every strength number on Page 1 is built up. Values are in virupas (60 virupas = 1 rupa).

\shot{stsb}
\shot{stbb}

Bhava Bala uses equal-house centres (ascendant degree plus 30° per house), because that is what \PL{} 9's
own Page 1 matches. The Yogi \& Bhava page uses Sripati houses by default.

\shot{sttime}
\shot{stav}

% =============================================================================
\section{Vargas}
\label{sec:vargas}
% =============================================================================
All the divisional charts \PL{} offers: the 16 Parashari charts (D1--D60), the Tajika charts and D150,
and \PL{}'s alternate versions of some vargas. The large chart shows one varga; click any small chart at
the bottom of the page to show it (the chosen one has a gold outline).

\shot{vamain}
\shot{vatable}

\begin{notebox}A planet exactly on an amsa boundary belongs to the amsa that begins there. The D11
``Iyer'' chart is a placeholder equal to the alternate D11 until Iyer's rule is supplied.\end{notebox}

% =============================================================================
\section{Dashas}
\label{sec:dashas}
% =============================================================================
The complete Vimshottari, down to the sixth level.

\shot{dacontrols}

\begin{itemize}[itemsep=2pt]
  \item \ui{Periods running on}: any date; \ui{Today} returns to today.
  \item \ui{Jump to the next period of}: a planet and a level; the tree opens at the next such period.
  \item \ui{Dasha year}: 365.2421875 days (\PL{}, default), 365.25 (Kala Chakra), 360, 327.85.
  \item \ui{Show times}: adds the clock time to every date (Prana and Deha always show it).
\end{itemize}

\shot{datree}
\shot{dapl}

\begin{qbox}\kq \emph{Why 365.2421875 when \PL{}'s option says 365.25?} Inside its program \PL{} uses
365.2421875 days. With it, all 10 dasha dates on the example's \PL{} 9 page match; with a true 365.25 only
5 of 10 do. A change here applies to every chart; \ui{Use the profile value} on the Settings page undoes it.
\end{qbox}

% =============================================================================
\section{KP (Research)}
\label{sec:kp}
% =============================================================================
This tab brings K.\,S. Krishnamurti's KP as his six \emph{KP Readers} teach it. Every rule carries its
book (B1--B6) and page, and was checked, where the books allow, on the Readers' own worked charts.

\shot{kpcalc}

\begin{tabularx}{\linewidth}{>{\bfseries}l X}
\toprule
Choice & What it does \\
\midrule
Readers mode (default) & What the Readers use: \PL{}'s \emph{Krishnamurti New} ayanamsa, the mean node, 365¼-day dasa years. \\
Automatic & \PL{}'s Krishnamurti ayanamsa, as on \PL{}'s KP pages (the Solar Fire profile keeps its Lahiri). \\
PL: Krishnamurti / New & Either of \PL{}'s two KP ayanamsas. \\
Same as the chart & The settings of every other page. \\
\bottomrule
\end{tabularx}

\subsection{Chart and tables}
\shot{kpchart}
\shot{kpcusps}
\shot{kpsig}

\subsection{Significators}
\shot{kprank}

Krishnamurti's order (B6 pp.\,94--95; B3 p.\,418): (1) planets in the stars of the occupants, (2) the
occupants, (3) planets in the stars of the owner, (4) the owner, (5) planets conjoined with or aspected
by these. A node acts for the planets it represents and comes just before them; a planet whose star holds
no planet is itself a strong significator. Ownership is by the cusp's sign only.

\shot{kppsets}

\subsection{Life matters}
\shot{kppromise}

For each of 14 matters (marriage, children, service, business, wealth, house, vehicle, education, abroad,
health, litigation, speculation, change of job, spiritual life) the cusp's sub lord decides: if it
signifies one of the matter's houses the matter is \emph{promised}; only the denial houses, \emph{denied};
both, \emph{mixed}. An asterisk marks a denial group derived by the 12th-from rule where the books print none.

\shot{kpmatter}

\begin{warnbox}\kpit Candidate periods are \textbf{candidates, not predictions}. On the Readers' own 11
marriage charts every wedding fell in such a period, but ranking by strength alone barely singled it out.
Tick \ui{Only periods whose lords are all ruling planets now}, then use \ui{Transit days}.\end{warnbox}

\subsection{Ruling planets and rectification}
Ruling planets lists the planets ruling the moment of judgement, now at the chart's place or at a
moment you set: the lords of the rising sign and star, of the Moon's star and sign, and of the day,
with the finer sub lords of the rising degree and the Moon below them. Rahu and Ketu enter as
agents of the planets they represent, and the list is ordered strongest first
(\autoref{fig:kprp}). Under it, the rectification assistant searches around the birth time for
moments whose rising sign, star, sub and sub-sub lords are all ruling planets
(\autoref{fig:kprect}). It proposes candidate times; checking them against events is yours to do.
\shot{kprp}
\shot{kprect}

\subsection{Transit timing}
Transit timing takes one period: the one running today, or the one whose \textbf{Transit days}
button you pressed in Life matters. It lists when a chosen body passes through the units whose
sign, star and sub lords are all lords of that period (\autoref{fig:kptr}). Tick extra planets
under \emph{Add ruling planets or significators} to widen the test. If the body makes no such
passage in the period, the app says so; try another body.
\shot{kptr}

\subsection{Horary 1--249}
\shot{kphor}
\shot{kphorc}

The three conditions (B6, KP6-061): (i) the cusp's sub lord is direct (if retrograde: delay until it
turns direct); (ii) its star lord is direct (if retrograde: denied); (iii) the sub lord signifies the
question's houses. The question list is B6's master table of house groups.

\subsection{KP Guidance}
\shot{kpguide}

\begin{keybox}\kidea What has been checked on the Readers' own examples: B6's significator tables; 45 horary
ascendants and 37 horary charts; the dasa periods B6 names at 32 dated events; 7 horary judgements; the
ruling planets of 15 judgement moments; 11 natal marriage charts; 5 transit dates; both rectification
cases. Where the books contradict themselves, the test names the case.\end{keybox}

% =============================================================================
\section{NTP (Research)}
\label{sec:ntp}
% =============================================================================
H.\,R. Seshadri Iyer's \emph{New Techniques of Prediction} applied to the chart. Only rules whose
conditions this chart meets are listed, each with its volume and page. Read them as the book's guidance,
not as a final judgement.

\shot{ntpguide}
\shot{ntpyoga}
\shot{ntpplanet}
\shot{ntpdasha}

\begin{notebox}The page uses the mean node, as Iyer says (``Take always Mean Rahu''); you can switch to
the chart's own setting. Houses are Sripati bhavas. Iyer's death and widowhood readings are deliberately
not shown. Marriage matching by Iyer's method is also on this tab (the same prospective match as on the
Matching tab).\end{notebox}

% =============================================================================
\section{Yogi \& Bhava}
\label{sec:yogi}
% =============================================================================
\shot{ybyogi}
\shot{ybbhava}
\shot{ybtable}

% =============================================================================
\section{Jaimini}
\label{sec:jaimini}
% =============================================================================
\shot{jakarakas}
\shot{jalagnas}
\shot{jaarudha}
\shot{jakm}
\shot{jaargala}
\shot{jadrishti}

% =============================================================================
\section{Transits}
\label{sec:transits}
% =============================================================================
\shot{trcontrols}
\shot{trpos}
\shot{trevents}

Rahu and Ketu are the true nodes unless changed in Settings; the true node wobbles, so it can turn direct
and retrograde again within a few days.

\shot{trsade}
\shot{trreturns}

% =============================================================================
\section{Matching}
\label{sec:matching}
% =============================================================================
Compatibility of two charts as \PL{} 7 reports it (Reports $\rightarrow$ Compatibility), plus Krishnamurti's
view. Say whether the first chart is the boy or the girl: the kootas count from the girl's Moon star to the
boy's. The example's partner below is a synthetic chart.

\shot{mentry}
\shot{mover}
\shot{mashta}
\shot{mmangal}
\shot{mdk}
\shot{mkp}

\begin{notebox}A real \PL{} 7.0.1 printout (Ashtakoot and Dash-Koota for one couple) is reproduced row for
row. \PL{} differs from some textbooks in two places: Tara is counted one way only, and two neutral sign
lords score 1 in Graha Maitri.\end{notebox}

% =============================================================================
\section{Settings}
\label{sec:settings}
% =============================================================================
\shot{secalc}

\begin{warnbox}\kpit A change here applies \textbf{to every chart in this browser} until you change it
back. The \ui{Calculation} profile on the chart form chooses all of these at once for one chart.\end{warnbox}

\shot{seplaces}

Saved places are kept in the browser on that device. Use \ui{Export places} on one device and
\ui{Import places…} on another to copy your list, and keep the file as a backup.

\shot{selook}

The round button at the top right switches the theme on any page, in the order Light, Dark, Warm night.
In every theme the chart is drawn as a dark plate with gold lines. On a phone the form wraps into one or
two columns, the panels sit in one column, the tabs scroll sideways and the Rasi chart shows planet names
only (the degrees are in the table under it).

% =============================================================================
\section{How far can I trust each number?}
\label{sec:trust}
% =============================================================================
``Checked against \PL{}'' means the app's number was compared with what \PL{} itself printed or showed.
Most checks use one \PL{} 9 screen (Page 1 of the example chart).

\begin{longtable}{>{\raggedright\arraybackslash}p{0.30\linewidth} >{\raggedright\arraybackslash}p{0.30\linewidth} >{\raggedright\arraybackslash}p{0.32\linewidth}}
\toprule
What you see & Checked against & Result \\
\midrule
\endhead
Planet longitudes (\PL{} 9) & \PL{} 9 Page 1 of the example & 9 of 9 equal to the displayed second \\
Ascendant & Same screen & Within 5″ (\PL{}'s atlas rounds coordinates) \\
Nakshatra, pada, KP lords & Same screen & All 10 points match \\
Vimshottari dates & Same screen & 10 of 10 (with 365.2421875-day years) \\
Page 1 strength fields & Same screen, 145 fields & No mismatches except Ketu's VB (open: \PL{} 14, app 9) \\
Divisional charts & Rules decoded from \PL{}'s program; PyJHora & Decoded and tested; not yet against a \PL{} 9 printout \\
Yogi point & \PL{}'s program; PyJHora & Matches; Avayogi and Duplicate Yogi rules open (selectable) \\
NTP & Iyer's worked examples & Exact, apart from 3 misprinted cells in the book \\
KP (Research) & The Readers' worked charts & As in \autoref{sec:kp}; book slips named in the tests \\
Matching & A \PL{} 7.0.1 printout & Row for row \\
Transits & Boundary crossings; \PL{} 7 transit records & Pass; not yet against a \PL{} 9 transit screen \\
Works offline & Automatic test, network off & Opens, casts, no network requests \\
\bottomrule
\end{longtable}

When a number in an open row matters for a judgement, check it in \PL{}. The written user guide
(\ex{docs/user\_guide/USER\_GUIDE.md}) has the complete table.

% =============================================================================
\section{Troubleshooting}
\label{sec:trouble}
% =============================================================================
\begin{description}[style=nextline, itemsep=4pt]
  \item[The green ``ready'' does not appear.] Wait a few seconds. If a red message starts \emph{The
        ephemeris could not start}, close the tab and open the file again, try another browser, or copy the
        file again (a partly copied file cannot start).
  \item[The planets are off by about an hour of motion.] The zone is wrong: check \ui{Zone} and
        \ui{Daylight saving}. India 1942--45 is 6.5 (or 5.5 with daylight saving ticked); before 1955 check
        which clock the time was written in; abroad, check summer time.
  \item[The date came out wrong.] Read the words under the Date box; the box may show month first.
  \item[My numbers differ from \PL{}'s.] Check the profile in the chart line; look for \ui{Use the profile
        value} in Settings; compare \ui{How was this computed?} with \PL{}; check \autoref{sec:trust}.
  \item[A \PL{} file will not open.] The app says \emph{This is not a Parashara's Light chart file} when the
        file has no birth data; choose the chart's own \ex{.xml} file.
  \item[Saved places have gone.] They live in this browser on this computer; a private window or clearing
        browser data starts empty. Keep an exported copy.
  \item[Printing a page.] Use the browser's Print (Ctrl+P or Cmd+P), Landscape for wide tables, and turn on
        \emph{Background graphics} so the chart plate prints with its dark background. The Light theme prints
        most clearly.
\end{description}

% =============================================================================
\section{Glossary}
\label{sec:glossary}
% =============================================================================
\begin{longtable}{>{\bfseries}p{0.24\linewidth} >{\raggedright\arraybackslash}p{0.70\linewidth}}
\toprule
Short & Meaning \\
\midrule
\endhead
As, Su … Ke & Ascendant (lagna), Sun, Moon, Mars, Mercury, Jupiter, Venus, Saturn, Rahu, Ketu \\
Ar … Pi & Aries, Taurus, Gemini, Cancer, Leo, Virgo, Libra, Scorpio, Sagittarius, Capricorn, Aquarius, Pisces \\
R, c & Retrograde, combust \\
SAV, AV & Sarvashtakavarga (all seven planets' bindus of a sign, total 337); ashtakavarga bindus of the planet in its sign \\
SB\%, SB\# & Shad Bala $\div$ required; number of Shad Bala groups at their minimum (0--5) \\
VB & Vimshopaka bala, out of 20 \\
Av3, Av5 & Three-state avastha (alert, dreaming, sleeping); five-state (infant, youth, adolescent, old, dead) \\
K7, AK … DK & Chara karaka (7-karaka scheme): Atma, Amatya, Bhratru, Matru, Pitru, (Putra), Gnati, Dara karaka \\
AL, A2 … A11, UL, KL & Arudha lagna, the other padas, upapada; Karakamsha lagna \\
D1 … D150 & Divisional charts (vargas) \\
BIN, STD, OPEN & Rule decoded from \PL{}'s program; textbook rule not yet confirmed; placeholder \\
Maha … Deha & The six levels of Vimshottari: Maha, Antar, Pratyantar, Sookshma, Prana, Deha \\
Sandhi, madhya & Beginning and centre of a bhava \\
KP; star / sub lords & Krishnamurti Paddhati; the lords of the nakshatra and of its Vimshottari subdivisions \\
RP & Ruling planets of a moment (KP) \\
CSL & Cusp sub lord (KP) \\
Ayanamsa & The difference between the tropical and sidereal zodiacs \\
$\Delta T$ (Delta-T) & The small difference between clock time and astronomical time \\
True / mean node & Rahu's exact (wobbling) position, or its smoothed average \\
Placidus & The house system \PL{} uses for cusps \\
Profile & A set of calculation choices: \PL{} 9, \PL{} 7, or Solar Fire / Kala Chakra \\
\bottomrule
\end{longtable}

% =============================================================================
\section*{About this document}
\addcontentsline{toc}{section}{About this document}
% =============================================================================
This guide was written for \app{} version 0.9. Every screenshot was taken automatically from the shipped
single file, and the numbered callouts were placed from the live positions of the elements on the page, so
the guide can be rebuilt exactly when the app changes (\ex{kaast2/media/pdf/}). It was prepared with the
help of an AI assistant (Claude, Anthropic) and checked against the app and its written user guide.

\end{document}
